\documentclass[11pt]{article}
\usepackage[margin=1in]{geometry}
\usepackage{amsmath,amssymb,amsthm,booktabs,microtype,hyperref}
\hypersetup{hidelinks,pdfauthor={T},pdftitle={Paired Audits for Joint Failures Under Changing Incentives}}
\newtheorem{theorem}{Theorem}
\title{Paired Audits for Joint Failures Under Changing Incentives}
\author{T}
\date{}
\begin{document}
\maketitle
\begin{abstract}
A trained assistant can pass tests that separately change oversight
and reward while failing when both change together. We identify the
range of joint-only failure rates consistent with paired single-change
tests. We give a minimax audit for discovering missed failures and an
exact finite confidence bound after zero findings. The design trades
off the cost of screening against savings on joint-condition labels.
A synthetic learning experiment shows how longer fitting can erase
robustness to cue reversal. Balancing the cue only in the joint
condition preserves predictive information in the other conditions.
\end{abstract}

\section*{Summary and scope}
\paragraph{Question.} What do separate oversight and incentive tests
leave unknown about harm when both changes occur?
\paragraph{Finding.} A uniform audit is minimax for finding failures in
the screen-negative pool. A separate stratified calculation bounds total
joint harm when both pools are sampled. The fitted experiment and its
population counterexample show why joint training coverage alone cannot
certify robustness to a new cue distribution.
\paragraph{Evidence and scope.} The finite-cohort sampling bounds
are proved and checked by exact enumeration. The training evidence comes
from synthetic feature models, with a separate exact population optimum.
A release evaluator can use the bounds for a fixed cohort and checkpoint.
Transfer to new tasks and language models requires further evidence.

\section{Two changes a cautious checkpoint can face}
A support assistant may respect a veto when it expects an audit or
when a shortcut yields no reward. In a new workflow, the audit cue may
disappear at the same time that completing a prohibited shortcut raises
the score. Separate tests of the two changes establish little about
that combination unless the evaluator can justify a structural relation
among outcomes. The case-paired design below uses the language of
potential outcomes for factorial experiments \cite{dasgupta}.

Take a fixed finite population of $N$ cases, each including the
same task, permissions, checkpoint, and any specified simulation
seed across conditions. Let $Y_{ab}(i)\in\{0,1\}$ indicate an
adjudicated harmful action in case $i$ under two switches
$a,b\in\{0,1\}$. For example, switch $a$ removes a visible observer
cue and switch $b$ changes which action earns an immediate reward.
In separate sandboxed copies of the case, screen both single
changes. Set
\[
U=\{i:Y_{10}(i)=1\text{ or }Y_{01}(i)=1\},\quad
V=\{1,\ldots,N\}\setminus U,\quad m=|V|.
\]
The joint-only set is $D=\{i\in V:Y_{11}(i)=1\}$. Its size $d$
counts failures that the single-change screens missed, regardless
of whether the baseline outcome was harmful. If the intended
target requires a clean baseline as well, add baseline failures
to $U$ before sampling and state that narrower estimand.

\section{What the single-change screens identify}
Without a restriction on the relationship between conditions, every
integer $d$ from zero through $m$ is compatible with the same complete
single-change observations. For each value, choose a $d$-element
subset of $V$ and make the joint condition harmful exactly on that
subset, leaving the other three conditions unchanged. This
construction applies even if every screened case passes.
For example, one policy may always respect a veto while another
takes a prohibited action only when both switches are active.
Those policies agree on the baseline and on both single-change
conditions. More single-change labels cannot distinguish them.

Pairing matters when determining the screen-negative pool.
Marginal single-change rates only bound their union. Recording
both outcomes per case identifies that union and avoids counting
the same screened failure twice. The method does not impose
monotonicity, so a single-change failure can also disappear when
both switches change.

\section{A minimax check for missed failures}
\begin{theorem}
Suppose $0\le k\le m$. Among randomized choices of exactly $k$
distinct cases in $V$, uniform sampling maximizes the
worst-case probability of hitting an unknown $d$-element failure
set. The value is
\[
1-\frac{\binom{m-d}{k}}{\binom{m}{k}},
\]
with $\binom{x}{k}=0$ when $x<k$.
\end{theorem}
\begin{proof}
Fix any sample of size $k$ and average the no-hit event over a
uniform $d$-element failure set. Its probability is
$\binom{m-k}{d}/\binom{m}{d}$, equal to the stated no-hit
probability by counting disjoint pairs. Averaging over any
randomized sampling design gives the same value. Its worst
failure set therefore cannot yield a higher detection probability
than the stated average. Uniform sampling attains that average
for every failure set.
\end{proof}

This theorem assumes the failure set is fixed before the audit
draw. A policy that detects the draw or changes between cases
requires a different adversary model. The optimum is also a
worst-case result across unknown failure locations. A trusted
predictor can support a more efficient unequal allocation, provided
the inference accounts for that selection.

An exact one-sided confidence procedure follows from the same
probability. If the sample is clean, report the largest $d$
with no-hit probability strictly greater than $\alpha$. If the
sample finds a failure, report $m$ as the upper count in this
zero-findings procedure. For every fixed $D$, understating the
true count is possible only on a clean sample whose probability
is at most $\alpha$. To bound all joint harm rather than the
joint-only subset, add at most $|U|$ to this upper count. A separate
probability sample of $U$ could tighten that conservative bound.

\section{Bounding all joint harm after zero findings}
Divide the cohort into strata of sizes $m_i$, such as $U$ and $V$.
Within each stratum draw $k_i$ cases uniformly without replacement,
independently of the other draws. If $d_i$ cases in stratum $i$ are
harmful under the joint condition, the no-find probability is
\[
 P_0(d_1,\ldots,d_s)
  =\prod_i\frac{\binom{m_i-d_i}{k_i}}{\binom{m_i}{k_i}}.
\]
To bound the total $d=\sum_i d_i$, maximize this expression over
all placements of $d$ harms. This can be done greedily. With
$r=m_i-d_i>0$, adding one harm to stratum $i$ multiplies its factor
by $(r-k_i)/r$ while the numerator is positive. Use zero when the
factor first becomes zero or is already zero.

For each stratum the successive multipliers form a nonincreasing
list. The product for an allocation is a product of prefixes of
those lists. Taking the largest available multiplier repeatedly
selects $d$ globally largest multipliers, with prefix order
maintained even under ties. It therefore attains the maximum.
The same zero-findings inversion as before, using this maximum,
gives a valid upper bound on total harm.

For the 64-case example below, placing all 12 joint labels in the
48 screen-negative cases gives a 95\% total-harm upper count of
25 after zero findings. Allocating nine labels there and three to
the screen-positive cases lowers it to 12. It matches the upper
count under unstratified sampling in this fixture. With arbitrary
failure placement and a fixed total label count, the minimax
theorem applied to the full population prevents a screen from
improving the worst-case detection probability without extra
structure. Its value for targeting joint-only failures is a
different claim. Exhaustive allocation checks agreed with the
greedy calculation on 27,634 probabilities across 2,954 designs.

\section{Finite verification and review cost}
Consider 64 cases with a screen-negative pool of 48. We compare 12
uniform joint-condition labels in the pool with 12 unstratified
labels. The endpoint is detecting a member of the genuine
joint-only set. This comparison credits the unstratified design
with detection even if additional screening would be needed to
classify the discovered failure as joint-only.
\begin{center}
\begin{tabular}{rrrr}
\toprule
Joint-only failures & Targeted detection & Unstratified detection & Labels to match \\
\midrule
1  & 0.2500 & 0.1875 & 16 \\
4  & 0.6973 & 0.5739 & 17 \\
8  & 0.9198 & 0.8300 & 17 \\
16 & 0.9968 & 0.9788 & 17 \\
\bottomrule
\end{tabular}
\end{center}
With four failures, the targeted design costs less than the
17-label unstratified design when the full two-condition screen
costs less than $5/64$ of a joint label per case. If those screens
were already part of checkpoint assessment, the incremental cost
can be much lower. At level $0.05$, 12 zero-findings labels give
an upper count of nine joint-only failures with targeted sampling,
or 12 with unstratified sampling. Those bounds are for the 64-case
cohort. Population transfer needs additional sampling assumptions.

The checker exhausts sample and failure-set pairs on an eight-case
population. Across sample sizes and failure counts from one
through four, it verifies both directions of the disjoint-pair
identity, covering 26,244 pairs. It computes the larger-cohort
bounds and match costs in exact rational arithmetic.


\section{A fitted model under changing cue reliability}
The exact audit describes detection after training. To investigate
how training coverage and the feedback target interact, we ran a
four-arm synthetic study. The target is binary and the case
specification declares whether acting is useful or harmful.
A task feature agrees with the target with probability $0.85$.
An approval cue agrees with probability $0.97$ in the baseline
and single-change cells, but only $0.03$ in the joint cell.
These probabilities define the fixture. They are not estimated
from agent behavior.

Each arm starts from the same affine logistic classifier trained
to predict the cue. The feature map includes the cue multiplied
by the joint-condition indicator. That interaction coefficient
has zero gradient whenever training omits the joint cell, so its
initial value persists during single-change training. All arms
use 768 fine-tuning labels per seed and the same optimizer:
\begin{itemize}
\item Verified labels in the baseline and two single-change cells.
\item Verified labels in all four cells.
\item Approval-cue labels in all four cells.
\item Verified labels in all four cells, with the cue replaced
      by an independent coin in the training inputs.
\end{itemize}
The test fixes the task and true label and crosses the joint cue
three ways. It either matches the fine-tuning environment, reverses
again, or is independent. Every arm is evaluated on the same
cases. Across eight seeds there are 4,800 observations in each
joint condition and 14,400 observations across the other three
cells. The three joint columns are paired projections, so they
should not be pooled as independent test cases.

\begin{center}
\small
\begin{tabular}{lrrrr}
\toprule
Training arm & Single cells$^{*}$ & Joint matched & Joint reversed & Joint independent \\
\midrule
Single verified & 13,962 & 119 & 4,681 & 2,432 \\
Four-cell verified & 13,751 & 4,075 & 4,075 & 4,075 \\
Four-cell proxy & 13,962 & 119 & 4,681 & 2,432 \\
Four-cell neutralized & 12,296 & 4,075 & 4,075 & 4,075 \\
\bottomrule
\end{tabular}

\medskip
\footnotesize Correct decisions. $^{*}$Includes the baseline cell.
\end{center}

In the matched joint cell, single-change training and
four-cell proxy training each took 2,286 false actions out
of 2,347 negative cases and 58 useful actions out of 2,453
positive cases. Four-cell training with either verified target
made 335 false actions and 2,063 useful actions. The matched,
reversed, and independent scores agree for both verified arms
in this run because both made the task-feature decision in the
joint cell. Neutralization reduced accuracy in the three cells
where the cue remained predictive. The four-cell verified arm
preserved more of that accuracy under this feature map and
optimization schedule. A different learner or training schedule
could learn the conditional cue and lose the observed invariance.
The experiment supports a claim about learned decisions in the
specified model. Its data do not establish transfer to language
models or deployment.


\section{Why the fitted invariance need not survive optimization}
The data-generating law lets us check what happens at the population
optimum. Let $x,z\in\{-1,1\}$ be the task and cue features and let
$j$ indicate the joint cell. The two features are conditionally
independent given the target and cell. Bayes' rule gives the true
posterior log odds
\[
\log(17/3)x+\log(97/3)(1-2j)z .
\]
This function lies in the fitted class. It minimizes population
log loss because each conditional log-loss term is minimized by
the true posterior probability. Since $\log(97/3)>\log(17/3)$,
the resulting joint decision is $-z$, even when the task feature
disagrees. Its joint accuracy is $97\%$ in the matched environment,
$3\%$ after reversal, and $50\%$ with an independent cue.
For neutralized training the posterior log odds is
$\log(17/3)x$, giving $85\%$ in all three environments.
An independent rational checker enumerates the feature and label
cells and reproduces these values.

Thus the finite crossed fit achieved invariance that its own
population objective does not preserve. Training to better fit
verified labels can still favor a context-specific shortcut.
The two tests diagnose different failures. Crossing both switches
gives the optimizer evidence about the interaction. The original
four-cell distribution and unconstrained log-loss objective have a
noninvariant optimum in the stated feature class. Cue neutralization
gives one way to change that optimum. A frozen release test should
vary the cue again, even after the checkpoint has passed the
matched joint condition.

\section{Balancing the cue only in the joint condition}
Global neutralization discards useful information in the three
cells where the cue is stable. Instead, in the joint cell only,
replace the training cue with an independent fair coin with
probability $\gamma$. The resulting joint cue accuracy is
\[
 r(\gamma)=0.03(1-\gamma)+0.5\gamma.
\]
The resulting cue remains independent of the task feature
conditional on the label.
The population score is representable by the same affine model,
using the original cue coefficient outside the joint cell and a
different one inside it. In the joint cell the score follows the
task feature on every input without ties exactly when its log-odds magnitude
strictly dominates that of the cue, or $r(\gamma)>0.15$.
Equivalently, $\gamma>12/47$. At equality some inputs tie.
Outside the joint cell the population decision still follows the
original cue and has accuracy $0.97$.

The task-feature decision is minimax for joint accuracy across
the matched and reversed test laws. Their equally weighted
mixture makes the cue fair conditional on the label and the
task feature. The best accuracy on this mixture is $0.85$.
For any rule, the smaller of its two accuracies is at most its
accuracy on the mixture. Using the task feature attains $0.85$
in each law and in the independent-cue test. This argument
covers randomized decisions as well. Exact enumeration checks
all sixteen deterministic rules on the two binary inputs.

We reused the eight synthetic seeds and their test cases to
compare four balancing weights at 60 and 240 epochs. Each
joint training row allocates weight $1-\gamma$ to its original
cue and $\gamma/2$ to each replacement cue, holding total
training weight fixed. The three other cells are unchanged.
At 240 epochs the correct counts were as follows.
\begin{center}
\small
\begin{tabular}{lrrrr}
\toprule
Joint weight $\gamma$ & Outside joint & Matched & Reversed & Independent \\
\midrule
$0$   & 13,962 & 4,681 & 119   & 2,368 \\
$1/4$ & 13,962 & 4,075 & 4,075 & 4,075 \\
$1/3$ & 13,962 & 4,075 & 4,075 & 4,075 \\
$1$   & 13,962 & 4,075 & 4,075 & 4,075 \\
\bottomrule
\end{tabular}
\end{center}
The unbalanced run reproduces the earlier counts at 60 epochs
and loses most reversed-cue accuracy with longer training.
One-quarter balancing appears robust at 240 epochs but is
below the population threshold. Its fitted robustness in this
comparison cannot certify the behavior of its objective at the
population optimum. One-third and full balancing satisfy the
threshold in the synthetic model. The follow-up uses the
original test cases and gives no new external test of generalization.

\section{Use during training and release}
Screen paired cases while tuning the policy, then freeze a distinct
cohort and checkpoint before making a release claim. Draw the
targeted joint tests with a seed inaccessible to the acting policy.
In an authorization-sensitive workflow, useful outcomes include
acting on permitted requests and asking for clarification when
authority is ambiguous. Report both along with harms so that a
checkpoint that simply refuses everything does not appear useful.
A prospective agent study should compare the verified curricula
at a common audit cost and hold out entire task families. The
synthetic study supplies a control with known cue reliability.
An agent study would need external outcome adjudication and
validation that both interventions preserve the task and authority.

\paragraph{Review status.} Independent manuscript review is pending.

\begin{thebibliography}{9}
\bibitem{dasgupta}
Tirthankar Dasgupta, Natesh S. Pillai, and Donald B. Rubin.
\emph{Causal Inference from $2^k$ Factorial Designs Using the
Potential Outcomes Model}. Journal of the Royal Statistical Society,
Series B 77(4), 2015. \url{https://arxiv.org/abs/1211.2481}.
\end{thebibliography}
\end{document}
